\documentclass[11pt]{article}
\usepackage[margin=1in]{geometry}
\usepackage{amsmath,amssymb}
\usepackage{graphicx}
\usepackage{booktabs}
\usepackage{hyperref}
\usepackage{natbib}

\title{Design Principles for Inflammasome Inhibition by Pyrin-Only-Proteins}
\author{Anonymous Authors}
\date{}

\begin{document}
\maketitle

\begin{abstract}
Inflammasomes are irreversible filamentous signaling platforms whose dysregulation drives autoinflammatory disease, cancer, and severe COVID-19. Pyrin-only-proteins (POPs) are endogenous regulators that inhibit pyrin domain (PYD) filament assembly, but their intrinsic target specificities and inhibitory mechanisms have remained speculative. Here we combine Rosetta-based interface energy analysis with quantitative biochemical polymerization assays and cell-based validation to define the target specificity and mechanism of POP1, POP2, and POP3. Contrary to the prevailing model, POP1 is a poor inhibitor of the central adaptor ASC; instead, it targets upstream receptor PYD filaments (AIM2, IFI16, NLRP3, NLRP6) through elongation inhibition. POP2 and POP3 both potently suppress ASC nucleation (prolonged lag phases in polymerization) while additionally inhibiting elongation of receptor filaments. Effective POP-PYD recognition requires a specific pattern of both favorable and unfavorable interactions at the six unique helical filament interface surfaces, as revealed by Rosetta energy scoring. POPs inhibit polymerization without co-assembling into filaments---in contrast to CARD-only-proteins (COPs)---and require super-stoichiometric concentrations for inhibition, especially in the presence of activating ligands. These findings revise the model of endogenous inflammasome regulation and reveal design principles for therapeutic targeting of specific inflammasome pathways.
\end{abstract}

\section{Introduction}

Inflammasomes are cytoplasmic multiprotein complexes that assemble into filamentous signaling platforms in response to pathogen-associated and damage-associated molecular patterns \citep{lu2014unified, stehlik2017review}. Assembly is initiated when sensor proteins (AIM2, IFI16, NLRP3, NLRP6, among others) oligomerize through their pyrin domains (PYDs), nucleating polymerization of the adaptor protein ASC through PYD-PYD interactions. ASC filaments then recruit and activate pro-caspase-1 through CARD-CARD interactions, leading to maturation of pro-inflammatory cytokines IL-1$\beta$ and IL-18 and pyroptotic cell death via gasdermin D pores \citep{sborgi2015gsdmd}.

A critical feature of inflammasome assembly is its irreversibility: once formed, PYD and CARD filaments do not dissociate and can self-perpetuate in a prion-like manner \citep{morrone2015assembly}. This creates a strong selective pressure for endogenous negative regulators. Pyrin-only-proteins (POPs) are small, single-domain PYD-like proteins that were identified as such regulators \citep{dejesus2007copsandpops}. The prevailing model, based primarily on overexpression studies, held that POP1 directly inhibits ASC \citep{khare2012pop1}. However, no systematic study has combined computational and biochemical approaches to define which PYD filaments each POP targets and through what mechanism.

\section{Methods}

\subsection{Rosetta Interface Analysis}

A honeycomb-like sideview of each PYD filament was constructed where each protomer provides six unique interface surfaces (Type 1a/b, 2a/b, 3a/b). The center protomer was replaced with each POP to compute interface energies ($\Delta G$ values in Rosetta Energy Units, REU) using Rosetta InterfaceAnalyzer \citep{rosettainterfaceanalyzer2015}. Cryo-EM filament structures of ASC$^{\text{PYD}}$ (PDB: 3j63), AIM2$^{\text{PYD}}$ (PDB: 7k3r), NLRP3$^{\text{PYD}}$ (PDB: 7pdz), and NLRP6$^{\text{PYD}}$ (PDB: 6ncv) served as templates.

\subsection{Biochemical Polymerization Assays}

FRET-based and fluorescence anisotropy (FA)-based quantitative polymerization assays measured assembly kinetics. MBP-tagged PYD constructs were triggered for assembly by TEV protease cleavage in the presence of increasing POP concentrations. Half-times ($t_{1/2}$) for polymerization were measured and IC$_{50}$ values calculated.

\subsection{Cell-Based Validation}

HEK293T cells were co-transfected with mCherry-tagged PYDs and eGFP-tagged POPs. Filament formation was quantified by automated fluorescence microscopy ($N \geq 4$ replicates).

\section{Results}

\subsection{POP1 Does Not Effectively Inhibit ASC}

Rosetta analysis predicted POP1 as a poor ASC$^{\text{PYD}}$ inhibitor, with multiple unfavorable interfaces. Cell-based assays confirmed this: POP1-eGFP at 600 and 1200 ng did not reduce ASC$^{\text{PYD}}$-mCherry filament counts, and FRET kinetics showed minimal lag prolongation. This directly contradicts the previous model.

\subsection{POP2 and POP3 Suppress ASC Nucleation}

Both POP2 and POP3 showed dose-dependent inhibition of ASC$^{\text{PYD}}$ filament formation in cells and markedly prolonged the lag phase of ASC polymerization in FRET assays, consistent with nucleation suppression. This was confirmed with full-length ASC.

\subsection{POP1 Targets Upstream Receptor PYD Filaments}

All three POPs inhibited AIM2$^{\text{PYD}}$ and IFI16$^{\text{PYD}}$ filament formation in cells and biochemical assays. POP1 and POP2 inhibited NLRP3$^{\text{PYD}}$ assembly, while POP3 showed minimal effect on NLRP3. All three POPs inhibited NLRP6$^{\text{PYD}}$. The primary mechanism for receptor PYD inhibition was elongation reduction rather than nucleation suppression.

\subsection{Design Principles from Rosetta Analysis}

Effective POP-PYD recognition required both favorable ($\Delta\Delta G \leq 3.5$ REU) and unfavorable ($\Delta\Delta G \geq 10.0$ REU) interactions at the six interface surfaces. Neither sequence similarity to the target PYD nor uniformly favorable interactions predicted inhibition effectiveness. Instead, the specific pattern of attractive and repulsive contacts determines specificity.

\subsection{POPs Do Not Co-Assemble into Filaments}

Negative-stain EM confirmed that POP1 (100 $\mu$M), POP2, and POP3 (20 $\mu$M after TEV cleavage) do not form filaments. This distinguishes their mechanism from COPs, which co-assemble into caspase-1 CARD filaments.

\subsection{Super-Stoichiometric Concentrations Required}

Effective inhibition required excess POP relative to PYD, especially in the presence of activating ligands such as dsDNA for AIM2 and IFI16. This is consistent with a model where POPs must compete with the assembly-promoting effect of the ligand.

\section{Discussion}

Our results fundamentally revise the model of POP-mediated inflammasome regulation. POP1 is not primarily an ASC inhibitor; rather, it targets upstream receptor PYD filaments. POP2 and POP3 serve as the primary ASC regulators through nucleation suppression. This revised specificity map---POP1 for receptors, POP2/POP3 for ASC plus additional receptor targets---provides a more nuanced view of endogenous inflammasome regulation.

The design principles revealed by Rosetta analysis have implications for therapeutic development. The requirement for both favorable and unfavorable interactions at specific interface surfaces suggests that effective inflammasome inhibitors must be designed to exploit the full hexavalent interaction geometry of PYD filaments, not simply maximize binding affinity at a single interface.

The super-stoichiometric requirement for POP inhibition provides an elegant regulatory logic: at baseline POP levels, inflammasome assembly can proceed when sufficiently stimulated, but elevated POP expression (e.g., during resolution of inflammation) can suppress further assembly.

\section{Conclusion}

We define the intrinsic target specificities and mechanisms of the three human POPs, revealing that POP1 targets upstream receptor PYDs while POP2 and POP3 suppress ASC nucleation. The design principles governing POP-PYD recognition---requiring specific patterns of favorable and unfavorable interactions across six interface surfaces---provide a framework for understanding endogenous inflammasome regulation and developing targeted therapeutic interventions.

\bibliographystyle{plainnat}
\bibliography{references}

\end{document}
